\section{Постановка задач}

Схема интерполяции в общем виде выглядит следующим образом:
\[
  f(t)\xrightarrow[\mathbf{R}]{}\{f_n\}^N_{n=0}\xrightarrow[\mathbf{I}]{}F(t)
\]

Здесь \( f(t) \) --- исходная функция, \( F(t) \) --- интерполирующая.

Как видно из схемы, на исходную функцию действует оператор сужения \( \mathbf{R} \) на сетку \( \{t_n\}^N_{n=0}\subset [a,b] \). Из неё формируется сеточная проекция функции \( f_n=\{f(t_n)\}_{n=0}^N \). Затем происходит неоднозначное восстановление исходной функции, которое определяется оператором интерполяции \( \mathbf{I} \).

В \textit{варианте №9} интерполяционные узлы определяются \textit{многочленами Чебышёва}, а оператор интерполяции --- \textit{интерполяционным многочленом в форме Лагранжа}.

То есть в ходе лабораторной работы будут выполнены следующие задачи:

\begin{enumerate}
    \item Разработка программного модуля для построения интерполяционного многочлена в форме Лагранжа $L_n(t)$.
    \item Использование корней многочлена Чебышёва в качестве узлов интерполяции на заданном отрезке $[a, b]$.
    \item Реализация итерационного процесса увеличения количества узлов $n$ до достижения заданной точности с критерием останова:
    $$|L_n(x) - L_{n-1}(x)| < \varepsilon, \text{ где } \varepsilon = 0.01$$
    \item Вычисление значения интерполирующей функции в заданной точке $x$ для различных входных функций и интервалов.
\end{enumerate}